% EGE 2027, task 7. Source IDs and URLs are kept in comments.
% Editable task bank. Append new BankTask blocks at the end of the document.
\documentclass[12pt,a4paper]{article}
\usepackage{physics-textbook}
\usepackage{graphicx,needspace}
\newif\ifShowSolutions
\ifdefined\TutorEdition
  \ShowSolutionstrue
\else
  \ShowSolutionsfalse
\fi
\newcounter{banktask}
\newcommand{\BankTask}[3]{%
  \refstepcounter{banktask}%
  \Needspace{35mm}\par\medskip\noindent\begin{minipage}{\linewidth}%
  \noindent\textbf{\thebanktask.}\quad #2\par
  \ifShowSolutions
    #3\par
    \ifdefined\TutorEdition
      \ifcsname TutorSource#1\endcsname
        \smallskip{\small\csname TutorSource#1\endcsname\par}%
      \fi
    \fi
  \fi
  \end{minipage}\par\medskip
}
\begin{document}
Справочные величины (если в условии не оговорено иное): $T/\text{К}=t/{}^\circ\text{C}+273$. $p_{\text{атм}}=10^5\,\text{Па}$. $R=8{,}31\,\text{Дж}/(\text{моль}\cdot\text{К})$.



% Source 0A3E4F: https://ege.fipi.ru/bank/questions.php?search=1&proj=BA1F39653304A5B041B656915DC36B38&qid=0A3E4F
\BankTask{0A3E4F}{
Горизонтальный цилиндрический сосуд с гладкими стенками разделён подвижным поршнем на две части. В одной части сосуда находится неон, в другой — аргон. Определите отношение средних кинетических энергий теплового движения молекул неона и аргона $\frac{E_{\text{н}}}{E_{\text{а}}}$, если поршень покоится, а отношение концентраций газов $\frac{n_{\text{н}}}{n_{\text{а}}}=\frac{1}{3}$.
}{
\smallskip\noindent\textbf{Решение от GPT}\par
При равновесии поршня давления газов равны: $p_{\text{н}}=p_{\text{а}}$. По основному уравнению МКТ $p=\frac{2}{3}nE$, поэтому $n_{\text{н}}E_{\text{н}}=n_{\text{а}}E_{\text{а}}$. Следовательно, $\frac{E_{\text{н}}}{E_{\text{а}}}=\frac{n_{\text{а}}}{n_{\text{н}}}=3$. Ответ: $3$.
}

% Source DC7C46: https://ege.fipi.ru/bank/questions.php?search=1&proj=BA1F39653304A5B041B656915DC36B38&qid=DC7C46
\BankTask{DC7C46}{
При уменьшении абсолютной температуры на $600\,\text{К}$ средняя кинетическая энергия теплового движения молекул аргона уменьшилась в $4$ раза. Какова конечная температура аргона (в $\text{К}$)?
}{
\smallskip\noindent\textbf{Решение от GPT}\par
Средняя кинетическая энергия $E=\frac{3}{2}kT$, поэтому $T_2=T_1/4$. Тогда $T_1-T_2=4T_2-T_2=600\,\text{К}$, откуда $T_2=200\,\text{К}$. Ответ: $200\,\text{К}$.
}

% Source 4C6FA1: https://ege.fipi.ru/bank/questions.php?search=1&proj=BA1F39653304A5B041B656915DC36B38&qid=4C6FA1
\BankTask{4C6FA1}{
Температура неона увеличилась с $27\,{}^{\circ}\text{C}$ до $177\,{}^{\circ}\text{C}$. Во сколько раз увеличилась средняя кинетическая энергия его молекул?
}{
\smallskip\noindent\textbf{Решение от GPT}\par
Абсолютные температуры: $T_1=(27+273)\,\text{К}=300\,\text{К}$, $T_2=(177+273)\,\text{К}=450\,\text{К}$. Поскольку $E=\frac{3}{2}kT$, получаем $\frac{E_2}{E_1}=\frac{450}{300}=1{,}5$. Ответ: в $1{,}5$ раза.
}

% Source 54C671: https://ege.fipi.ru/bank/questions.php?search=1&proj=BA1F39653304A5B041B656915DC36B38&qid=54C671
\BankTask{54C671}{
Во сколько раз уменьшится абсолютная температура разреженного одноатомного газа, если среднеквадратичная скорость теплового движения его молекул уменьшится в \(5\) раз?
}{
\smallskip\noindent\textbf{Решение от GPT}\par
Из \(\frac{m v_{\text{кв}}^2}{2}=\frac{3}{2}kT\) следует \(T\propto v_{\text{кв}}^2\). Поэтому \(\frac{T_1}{T_2}=\left(\frac{v_{\text{кв},1}}{v_{\text{кв},2}}\right)^2=5^2=25\).

Ответ: температура уменьшится в \(25\) раз.
}

% Source 9D4F49: https://ege.fipi.ru/bank/questions.php?search=1&proj=BA1F39653304A5B041B656915DC36B38&qid=9D4F49
\BankTask{9D4F49}{
Первоначальное давление идеального одноатомного газа в сосуде равнялось \(180\,\text{кПа}\). Объём сосуда увеличили, при этом концентрация молекул газа уменьшилась в \(3\) раза. Одновременно сосуд нагрели так, что средняя кинетическая энергия теплового движения молекул газа увеличилась в \(2\) раза. Каким в результате этого стало давление \(p\) газа в сосуде (в $\text{кПа}$)?
}{
\smallskip\noindent\textbf{Решение от GPT}\par
По основному уравнению МКТ \(p=\frac{2}{3}n\overline{E}_{\text{к}}\). Тогда \(\frac{p_2}{p_1}=\frac{n_2}{n_1}\frac{\overline{E}_{\text{к},2}}{\overline{E}_{\text{к},1}}=\frac{2}{3}\), откуда \(p_2=180\cdot\frac{2}{3}=120\,\text{кПа}\).

Ответ: \(120\,\text{кПа}\).
}

% Source 3AB2B2: https://ege.fipi.ru/bank/questions.php?search=1&proj=BA1F39653304A5B041B656915DC36B38&qid=3AB2B2
\BankTask{3AB2B2}{
В ходе эксперимента давление разреженного газа в сосуде снизилось в \(5\) раз, а средняя энергия теплового движения его молекул уменьшилась в \(2\) раза. Во сколько раз уменьшилась при этом концентрация молекул газа в сосуде?
}{
\smallskip\noindent\textbf{Решение от GPT}\par
Из \(p=\frac{2}{3}n\overline{E}_{\text{к}}\) получаем \(\frac{n_2}{n_1}=\frac{p_2}{p_1}\frac{\overline{E}_{\text{к},1}}{\overline{E}_{\text{к},2}}=\frac{1}{5}\cdot2=\frac{2}{5}\). Следовательно, \(\frac{n_1}{n_2}=\frac{5}{2}=2{,}5\).

Ответ: концентрация уменьшилась в \(2{,}5\) раза.
}

% Source 772AAE: https://ege.fipi.ru/bank/questions.php?search=1&proj=BA1F39653304A5B041B656915DC36B38&qid=772AAE
\BankTask{772AAE}{
В сосуде содержится разреженный аргон, абсолютная температура которого равна \(150\ \text{К}\). Концентрацию аргона уменьшили в \(2\) раза, при этом его давление увеличилось в \(3\) раза. Определите абсолютную температуру газа в конечном равновесном состоянии. (в $\text{К}$).
}{
\smallskip\noindent\textbf{Решение от GPT}\par
Из уравнения \(p=nkT\) следует
\[
\frac{T_2}{T_1}=\frac{p_2}{p_1}\frac{n_1}{n_2}=3\cdot2=6.
\]
Следовательно, \(T_2=6\cdot150=900\ \text{К}\).
Ответ: \(900\ \text{К}\).
}

% Source AD774B: https://ege.fipi.ru/bank/questions.php?search=1&proj=BA1F39653304A5B041B656915DC36B38&qid=AD774B
\BankTask{AD774B}{
В сосуде неизменного объёма находится идеальный газ. Во сколько раз нужно увеличить количество газа в сосуде, чтобы после уменьшения абсолютной температуры газа в \(2\) раза его давление стало вдвое больше начального?
}{
\smallskip\noindent\textbf{Решение от GPT}\par
По уравнению \(pV=\nu RT\) при неизменном объёме
\[
\frac{\nu_2}{\nu_1}=\frac{p_2}{p_1}\frac{T_1}{T_2}=2\cdot2=4.
\]
Ответ: в \(4\) раза.
}

% Source 8E16FE: https://ege.fipi.ru/bank/questions.php?search=1&proj=BA1F39653304A5B041B656915DC36B38&qid=8E16FE
\BankTask{8E16FE}{
Объём \(1\ \text{моль}\) водорода в сосуде при температуре \(T_0\) и давлении \(p_0\) равен \(20\ \text{л}\). Каков объём \(3\ \text{моль}\) водорода при том же давлении и вдвое большей температуре (в $\text{л}$)?
}{
\smallskip\noindent\textbf{Решение от GPT}\par
Из уравнения \(pV=\nu RT\) при одинаковом давлении получаем
\[
\frac{V_2}{V_1}=\frac{\nu_2}{\nu_1}\frac{T_2}{T_1}=3\cdot2=6.
\]
Тогда \(V_2=6\cdot20=120\ \text{л}\).
Ответ: \(120\ \text{л}\).
}

% Source AB117F: https://ege.fipi.ru/bank/questions.php?search=1&proj=BA1F39653304A5B041B656915DC36B38&qid=AB117F
\BankTask{AB117F}{
С идеальным газом провели изотермический процесс, в котором в результате увеличения объёма газа на \(300\ \text{дм}^{3}\) его давление уменьшилось в \(2\) раза.

Масса газа постоянна. Каким был первоначальный объём газа (в $\text{дм}^3$)?
}{
\smallskip\noindent\textbf{Решение от GPT}\par
По закону Бойля — Мариотта \(p_1V_1=p_2V_2\). Так как \(p_2=p_1/2\), то \(V_2=2V_1\). Следовательно, \(V_2-V_1=V_1=300\ \text{дм}^{3}\).

Ответ: \(300\ \text{дм}^{3}\).
}

% Source 3845F2: https://ege.fipi.ru/bank/questions.php?search=1&proj=BA1F39653304A5B041B656915DC36B38&qid=3845F2
\BankTask{3845F2}{
Разреженный неон в количестве \(5\ \text{моль}\) изобарно расширяется. Во сколько раз увеличится объём газа при повышении его абсолютной температуры в \(3\) раза?
}{
\smallskip\noindent\textbf{Решение от GPT}\par
При постоянных давлении и количестве вещества \(V/T=\mathrm{const}\), поэтому \(V_2/V_1=T_2/T_1=3\).

Ответ: в \(3\) раза.
}

% Source 349197: https://ege.fipi.ru/bank/questions.php?search=1&proj=BA1F39653304A5B041B656915DC36B38&qid=349197
\BankTask{349197}{
Идеальный газ в количестве \(5\ \text{моль}\) нагревается от \(t_1=127\,^{\circ}\text{С}\) до \(t_2=227\,^{\circ}\text{С}\) в сосуде постоянного объёма. Чему равно отношение давлений \(\frac{p_1}{p_2}\) идеального газа?
}{
\smallskip\noindent\textbf{Решение от GPT}\par
При постоянных объёме и количестве вещества \(p/T=\mathrm{const}\). Температуры: \(T_1=127+273=400\ \text{К}\), \(T_2=227+273=500\ \text{К}\). Тогда
\[\frac{p_1}{p_2}=\frac{T_1}{T_2}=\frac{400}{500}=0{,}8.\]
Ответ: \(0{,}8\).
}

% Source 5F4252: https://ege.fipi.ru/bank/questions.php?search=1&proj=BA1F39653304A5B041B656915DC36B38&qid=5F4252
\BankTask{5F4252}{
\begin{center}
\begin{tikzpicture}[x=1cm,y=0.065cm,>=stealth]
\draw[->] (0,0)--(4.7,0) node[right] {$V,\,\text{м}^{3}$};
\draw[->] (0,0)--(0,59) node[left] {$p,\,\text{кПа}$};
\node[below left] at (0,0) {$0$};
\foreach \x in {1,2,3} {\draw (\x,-2)--(\x,2);}
\node[below] at (1,-2) {$1$};
\node[below] at (3,-2) {$3$};
\foreach \y in {20,40} {\draw (-0.12,\y)--(1.25,\y);\node[left] at (-0.12,\y) {$\y$};}
\draw (2.5,20)--(3,20);
\draw (2.5,40)--(3,40);
\draw[thick] (3,20)--(3,40);
\fill (3,20) circle (1.5pt) node[right] {$1$};
\fill (3,40) circle (1.5pt) node[right] {$2$};
\end{tikzpicture}
\end{center}
На $pV$-диаграмме представлен процесс $1$--$2$, в котором участвует разреженный аргон ($p$ — давление аргона, $V$ — его объём). Абсолютная температура газа в состоянии $1$ равна $250\,\text{К}$. Какая абсолютная температура соответствует состоянию $2$, если масса газа остаётся неизменной (в $\text{К}$)?
}{
\smallskip\noindent\textbf{Решение от GPT}\par
При неизменной массе газа $pV=\nu RT$. Объём постоянен, поэтому
\[
T_2=T_1\frac{p_2}{p_1}=250\cdot\frac{40}{20}=500\,\text{К}.
\]
Ответ: $500\,\text{К}$.
}

% Source B86966: https://ege.fipi.ru/bank/questions.php?search=1&proj=BA1F39653304A5B041B656915DC36B38&qid=B86966
\BankTask{B86966}{
\begin{center}
\begin{tikzpicture}[x=0.012cm,y=0.8cm,>=stealth]
\draw[->] (0,0)--(530,0) node[right] {$T,\,\text{К}$};
\draw[->] (0,0)--(0,5) node[left] {$V,\,\text{м}^{3}$};
\node[below left] at (0,0) {$0$};
\draw (273,-0.08)--(273,0.08);
\node[below] at (273,-0.08) {$273$};
\draw (35,1)--(210,1);
\draw (65,4)--(245,4);
\node[left] at (0,1) {$1$};
\node[left] at (0,4) {$4$};
\draw[dashed] (273,0)--(273,1);
\draw (273,1)--(273,4);
\draw[->,thick] (273,1)--(273,2.5);
\fill (273,1) circle (1.3pt) node[right] {$1$};
\fill (273,4) circle (1.3pt) node[right] {$2$};
\end{tikzpicture}
\end{center}
На рисунке представлен график зависимости объёма идеального газа от его температуры в некотором процессе. В состоянии $1$ давление газа было равно нормальному атмосферному давлению. Какое давление соответствует состоянию $2$, если масса газа остаётся неизменной (в $\text{кПа}$)?
}{
\smallskip\noindent\textbf{Решение от GPT}\par
По графику $T_1=T_2=273\,\text{К}$, $V_1=1\,\text{м}^{3}$, $V_2=4\,\text{м}^{3}$. Поэтому $p_1V_1=p_2V_2$. Принимая нормальное атмосферное давление $p_1\approx100\,\text{кПа}$, получаем
\[
p_2=p_1\frac{V_1}{V_2}=100\cdot\frac{1}{4}=25\,\text{кПа}.
\]
Ответ: $25\,\text{кПа}$.
}

% Source 79B504: https://ege.fipi.ru/bank/questions.php?search=1&proj=BA1F39653304A5B041B656915DC36B38&qid=79B504
\BankTask{79B504}{
\begin{center}
\begin{tikzpicture}[x=0.11cm,y=5.5cm,>=stealth]
\draw[thin,gray] (0,0) grid[xstep=10,ystep=0.2] (40,0.6);
\draw[->] (0,0)--(51,0) node[right] {$T,\,\text{К}$};
\draw[->] (0,0)--(0,0.82) node[left] {$p,\,10^{4}\,\text{Па}$};
\foreach \x/\lab in {0/290,10/300,20/310,30/320,40/330} {
\draw (\x,0)--(\x,-0.015);
\node[below] at (\x,-0.015) {$\lab$};
}
\foreach \y/\lab in {0/{1{,}8},0.2/{2{,}0},0.4/{2{,}2},0.6/{2{,}4}} {
\draw (-0.7,\y)--(0,\y);
\node[left] at (-0.7,\y) {$\lab$};
}
\draw[thick] (0,0.1333)--(44,0.4267);
\end{tikzpicture}
\end{center}
На рисунке показан график зависимости давления газа в запаянном сосуде от его температуры. Объём сосуда равен $0{,}25\,\text{м}^{3}$. Какое приблизительно количество газообразного вещества содержится в этом сосуде (в $\text{моль}$, с округлением до целого)?
}{
\smallskip\noindent\textbf{Решение от GPT}\par
По графику при $T=300\,\text{К}$ давление $p\approx2{,}0\cdot10^{4}\,\text{Па}$. Из уравнения идеального газа
\[
\nu=\frac{pV}{RT}=\frac{2{,}0\cdot10^{4}\cdot0{,}25}{8{,}31\cdot300}\approx2{,}01\,\text{моль}.
\]
После округления ответ: $2\,\text{моль}$.
}

% Source 90F3B1: https://ege.fipi.ru/bank/questions.php?search=1&proj=BA1F39653304A5B041B656915DC36B38&qid=90F3B1
\BankTask{90F3B1}{
При изобарном расширении постоянной массы гелия при нормальном атмосферном давлении его объём увеличился на $0{,}2\,\text{м}^3$. Определите изменение внутренней энергии гелия. (в $\text{кДж}$).
}{
\smallskip\noindent\textbf{Решение от GPT}\par
Для одноатомного идеального газа $U=\frac32pV$. При постоянном давлении $\Delta U=\frac32p\Delta V=\frac32\cdot10^5\cdot0{,}2=3\cdot10^4\,\text{Дж}=30\,\text{кДж}$. Ответ: $30\,\text{кДж}$.
}

% Source A1DD9C: https://ege.fipi.ru/bank/questions.php?search=1&proj=BA1F39653304A5B041B656915DC36B38&qid=A1DD9C
\BankTask{A1DD9C}{
Из металлического баллона, в котором находилось $3\,\text{моль}$ идеального газа под давлением $3\cdot10^5\,\text{Па}$, медленно, без изменения температуры выпустили $0{,}6\,\text{моль}$ газа. Определите давление газа в конечном состоянии. (в $\text{кПа}$).
}{
\smallskip\noindent\textbf{Решение от GPT}\par
При постоянных $T$ и $V$ давление пропорционально количеству вещества: $p_2=p_1\nu_2/\nu_1=3\cdot10^5(3-0{,}6)/3=2{,}4\cdot10^5\,\text{Па}$. Ответ: $240\,\text{кПа}$.
}

% Source 67AF6A: https://ege.fipi.ru/bank/questions.php?search=1&proj=BA1F39653304A5B041B656915DC36B38&qid=67AF6A
\BankTask{67AF6A}{
С идеальным газом происходит изохорный процесс, в котором в результате увеличения абсолютной температуры газа в $2$ раза его давление возросло на $75\,\text{кПа}$. Масса газа постоянна. Каково было первоначальное давление газа (в $\text{кПа}$)?
}{
\smallskip\noindent\textbf{Решение от GPT}\par
При постоянных массе и объёме $p/T=\mathrm{const}$, поэтому $p_2=2p_1$. Следовательно, $p_2-p_1=p_1=75\,\text{кПа}$. Ответ: $75\,\text{кПа}$.
}

% Source EE7CA0: https://ege.fipi.ru/bank/questions.php?search=1&proj=BA1F39653304A5B041B656915DC36B38&qid=EE7CA0
\BankTask{EE7CA0}{
Идеальный одноатомный газ находится в сосуде с жёсткими стенками объёмом \(0{,}6\ \text{м}^{3}\). При охлаждении его давление снизилось на \(3\cdot10^{3}\ \text{Па}\).

На сколько уменьшилась внутренняя энергия газа (в $\text{кДж}$)?
}{
\smallskip\noindent\textbf{Решение от GPT}\par
Для одноатомного идеального газа \(U=\frac32pV\). При постоянном объёме уменьшение внутренней энергии равно
\[
U_1-U_2=\frac32V(p_1-p_2)=\frac32\cdot0{,}6\cdot3\cdot10^3=2700\ \text{Дж}=2{,}7\ \text{кДж}.
\]
Ответ: \(2{,}7\ \text{кДж}\).
}

% Source 1B1B5D: https://ege.fipi.ru/bank/questions.php?search=1&proj=BA1F39653304A5B041B656915DC36B38&qid=1B1B5D
\BankTask{1B1B5D}{
При изобарном расширении неона его объём увеличился на \(0{,}4\,\text{м}^3\), а внутренняя энергия увеличилась на \(6\,\text{кДж}\). Определите давление неона. (в $\text{кПа}$).
}{
\smallskip\noindent\textbf{Решение от GPT}\par
Для одноатомного идеального газа \(U=\frac32pV\). При постоянном давлении \(\Delta U=\frac32p\Delta V\), откуда
\[
p=\frac{2\Delta U}{3\Delta V}=\frac{2\cdot6\cdot10^3}{3\cdot0{,}4}=10^4\,\text{Па}=10\,\text{кПа}.
\]
Ответ: \(10\,\text{кПа}\).
}

\end{document}
